\section{Objetivos del proyecto}

\subsection{Objetivo general}
\begin{indentedsubsection}
De acuerdo con el diagnóstico y el Plan de Implementación, el objetivo principal es establecer una ruta estratégica y operativa que guíe la adopción del sistema de gestión en la nube basado en Odoo para Perú Flores Travel Agency Tours Operador E.I.R.L.. El proyecto busca erradicar la fragmentación operativa actual ---basada en cuadernos físicos, hojas de Excel y mensajes de WhatsApp--- mediante una plataforma centralizada que agilice la gestión comercial y automatice el cálculo y envío de cotizaciones de paquetes turísticos.

Este objetivo se centra en devolverle a la gerencia el control de su tiempo: pasar de invertir horas en sumas manuales y transcripciones de datos, a administrar un embudo de ventas ágil, profesionalizando la presentación de propuestas hacia el turista extranjero y blindando la información de la agencia.
\end{indentedsubsection}

\subsection{Objetivos específicos}
\begin{indentedsubsection}
Para materializar el objetivo general, se han trazado las siguientes metas puntuales basadas en las funcionalidades estándar (out-of-the-box) de la plataforma:

\begin{itemize}
    \item \textbf{a. Implementar un CRM y base de datos centralizada en la nube:} desplegar el módulo CRM para unificar el registro de pasajeros (PAX) y el directorio de proveedores (guías, transporte, hoteles), eliminando la dependencia de libretas físicas. Esto permitirá organizar a los clientes en un embudo de ventas claro: Nuevo, Calificado, Propuesta y Ganado.
    \item \textbf{b. Automatizar el flujo de cotizaciones y armado de paquetes:} configurar un cotizador digital con un catálogo de servicios tarifados que reemplace el copiado y pegado en Word y las sumas en Excel. Esto permitirá calcular prorrateos, simular porcentajes de ganancia dinámicos y generar instantáneamente propuestas comerciales en formato PDF.
    \item \textbf{c. Gestionar eficientemente cobros y pagos fraccionados:} parametrizar el registro de pagos dentro del sistema para llevar un control exacto de los adelantos (``a cuenta'') y los saldos pendientes, adaptándose a la realidad de la agencia que gestiona cobros variables mediante Western Union o depósitos bancarios.
    \item \textbf{d. Sistematizar la retroalimentación y calidad post-servicio:} diseñar e implementar el módulo de Encuestas (Feedback) para enviar automáticamente formularios de satisfacción a los turistas al finalizar su itinerario. El objetivo es consolidar métricas de calidad de los proveedores y recopilar testimonios que potencien el ``boca a boca'' digital.
    \item \textbf{e. Desarrollar competencias digitales operativas:} ejecutar un plan de capacitación intensiva y práctica (``aprender haciendo'') dirigido a la gerencia y a su equipo de apoyo familiar. El enfoque garantizará que los usuarios sean 100\% autónomos en la creación de paquetes turísticos, generación de PDFs y gestión del CRM en su día a día.
\end{itemize}
\end{indentedsubsection}
